% part: intuitionistic-logic
% chapter: soundness-completeness
% section: lindenbaum

\documentclass[../../../include/open-logic-section]{subfiles}

\begin{document}

\olfileid{int}{sc}{lin}

\olsection{Lindenbaum 引理}

证明直觉主义逻辑的完备性定理，需要假设 $\Gamma \Proves/ !A$ 并构造一个满足 $\Gamma$ 但不满足 $!A$ 的模型 $\mSat{M}{\Gamma}$。

在经典逻辑中，可推导关系可以化归为一致性的概念，因为公式 $!A$ 可从一公式集推导出，当且仅当该集合与 $!A$ 的否定合在一起是不一致的。这在直觉主义逻辑中是不可能的。在直觉主义逻辑中，如果 $\lnot!A$ 不一致，我们只能得到 $\Proves \lnot\lnot !A$。由于 $\lnot\lnot!A \lif !A$ 在直觉主义逻辑中一般不成立，我们无法得出 $\Proves !A$ 的结论。

因此，在构造模型 $\mModel{M}$ 时，我们需要保留公式 $!A$ 不可推导这一性质，从而将无法使用完全集 $\Gamma^* \supseteq \Gamma$ 来构建模型 $\mModel{M}$，因为在每一个完全集 $\Gamma^*$ 中，都有 $\Gamma^* \Proves !A \lor \lnot !A$。

我们将不使用完全集 $\Gamma^*$，而是使用公式素集的概念：

\begin{defn}\ollabel{defn:prime}
  公式集合 $\Gamma$ 是\emph{素的}，当且仅当
  \begin{enumerate}
  \item\ollabel{defn:prime1} $\Gamma$ 是一致的，即 $\Gamma \Proves/ \lfalse$；
  \item\ollabel{defn:prime2} 若 $\Gamma \Proves !A$，则 $!A \in \Gamma$；且
  \item\ollabel{defn:prime3} 若 $!A \lor !B \in \Gamma$，则 $!A \in \Gamma$ 或 $!B \in \Gamma$。
  \end{enumerate}
\end{defn}

\begin{lem}[Lindenbaum 引理]
  \ollabel{lem:lindenbaum} 如果 $\Gamma \Proves/ !A$，那么存在一个 $\Gamma^* \supseteq \Gamma$，使得 $\Gamma^*$ 是素的并且 $\Gamma^* \Proves/ !A$。
\end{lem}

\begin{proof}
  令 $!B_1 \lor !C_1$, $!B_2 \lor !C_2$, \dots{} 为所有形如 $!B \lor !C$ 的公式的一个枚举。
  我们将定义一个递增的公式集序列 $\Gamma_n$，其中每个 $\Gamma_{n+1}$ 定义为 $\Gamma_n$ 加上一个新公式。
  $\Gamma^*$ 将是所有 $\Gamma_n$ 的并集。
  新公式的选取是为了确保 $\Gamma^*$ 是素的并且仍满足 $\Gamma^* \Proves/ !A$。
  这意味着在每一步，我们应当找到第一个析取式 $!B_i \lor !C_i$，使得：
  \begin{enumerate}
  \item\ollabel{gamma-1} $\Gamma_n \Proves !B_i \lor !C_i$
  \item\ollabel{gamma-2} $!B_i \notin \Gamma_n$ 且 $!C_i \notin \Gamma_n$
  \end{enumerate}
  我们将把 $!B_i$ 或 $!C_i$ 加到 $\Gamma_n$ 中：若 $\Gamma_n \cup \{!B_i\}
  \Proves/ !A$ 则加 $!B_i$，否则加 $!C_i$。我们需要证明这是可行的。目前，我们先定义 $i(n)$ 为满足条件 \olref{gamma-1} 和 \olref{gamma-2} 的最小 $i$。

  定义 $\Gamma_0 = \Gamma$ 且
  \[
  \Gamma_{n+1} = \begin{cases}
    \Gamma_n \cup \{!B_{i(n)}\} &
    \text{if $\Gamma_n \cup \{!B_{i(n)}\} \Proves/ !A$} \\
    \Gamma_n \cup \{!C_{i(n)}\} & \text{否则}
    \end{cases}
  \]
  如果 $i(n)$ 无定义，即只要 $\Gamma_n \Proves !B \lor !C$，就有 $!B \in \Gamma_n$ 或 $!C \in \Gamma_n$，则令 $\Gamma_{n+1}
  = \Gamma_n$。现在令 $\Gamma^* = \bigcup_{n=0}^\infty \Gamma_n$。

  首先我们证明，对所有 $n$，都有 $\Gamma_n \Proves/ !A$。我们对 $n$ 进行归纳。
  当 $n = 0$ 时，断言由引理的假设成立，即 $\Gamma \Proves/ !A$。
  若 $n>0$，我们需要证明：如果 $\Gamma_n \Proves/ !A$，那么 $\Gamma_{n+1} \Proves/ !A$。
  如果 $i(n)$ 无定义，则 $\Gamma_{n+1} = \Gamma_n$，无需证明。
  所以假设 $i(n)$ 有定义。为简单起见，令 $i = i(n)$。

  我们将证明该断言的逆否命题。假设 $\Gamma_{n+1} \Proves !A$。
  根据构造，$\Gamma_{n+1} = \Gamma_n \cup
  \{!B_i\}$ 若 $\Gamma_n \cup \{!B_i\} \Proves/ !A$，否则 $\Gamma_{n+1} = \Gamma_n \cup \{!C_i\}$。
  显然不能是第一种情况，因为那样会 $\Gamma_{n+1} \Proves/ !A$。
  因此，$\Gamma_n \cup
  \{!B_i\} \Proves !A$ 且 $\Gamma_{n+1} = \Gamma_n \cup \{!C_i\}$。
  根据 $i(n)$ 的定义，我们有 $\Gamma _n \Proves !B_i \lor
  !C_i$。
  我们有 $\Gamma_n \cup \{!B_i\} \Proves !A$。我们还有 $\Gamma_{n+1} = \Gamma_n \cup \{!C_i\} \Proves !A$。
  因此，$\Gamma_n \Proves !A$，这正是我们想要证明的。

  如果 $\Gamma^* \Proves !A$，那么会存在某个有限子集 $\Gamma' \subseteq \Gamma^*$ 使得 $\Gamma' \Proves !A$。
  每个 $!D \in \Gamma'$ 必定在某个 $i$ 对应的 $\Gamma_i$ 中。
  令 $n$ 为这些 $i$ 中最大的那个。
  因为如果 $i \le n$ 则有 $\Gamma_i \subseteq \Gamma_n$，所以 $\Gamma' \subseteq \Gamma_n$。
  但这样一来 $\Gamma_n \Proves !A$，就与我们上面证明的 $\Gamma_n \Proves/ !A$ 矛盾了。

  最后，我们证明 $\Gamma^*$ 是素的，即满足 \olref{defn:prime} 中的条件 \olref{defn:prime1}、\olref{defn:prime2} 和 \olref{defn:prime3}。

  首先，由于 $\Gamma^* \Proves/ !A$，$\Gamma^*$ 是一致的，故 \olref{defn:prime1} 成立。

  我们现在证明：如果 $\Gamma^* \Proves !B \lor !C$，那么要么 $!B \in \Gamma^*$，要么 $!C \in \Gamma^*$。
  这就证明了 \olref{defn:prime3}，因为如果 $!B \lor !C \in \Gamma^*$，那么也有 $\Gamma^* \Proves !B \lor !C$。
  为此，假设 $\Gamma^* \Proves !B \lor !C$ 但 $!B \notin \Gamma^*$ 且 $!C \notin \Gamma^*$。
  由于 $\Gamma^* \Proves !B \lor !C$，存在某个 $n$ 使得 $\Gamma_n \Proves !B \lor !C$。
  $!B \lor !C$ 出现在所有析取式的枚举中，设其为 $!B_j \lor !C_j$。
  $!B_j \lor !C_j$ 满足 $i(n)$ 定义中的条件，即我们有 $\Gamma_n \Proves !B_j
  \lor !C_j$，且 $!B_j \notin \Gamma_n$ 与 $!C_j \notin \Gamma_n$。
  在每一阶段，满足条件的析取式 $!B_i \lor !C_i$ 至少减少一个（因为在每一阶段我们添加了 $!B_i$ 或 $!C_i$），所以在某个阶段 $m$ 我们会有 $j = i(m)$。
  但那样的话，要么 $!B \in \Gamma_{m+1}$，要么 $!C \in \Gamma_{m+1}$，这与 $!B \notin \Gamma^*$ 且 $!C \notin \Gamma^*$ 的假设矛盾。

  现在假设 $\Gamma^* \Proves !B$。那么 $\Gamma^* \Proves !B \lor !B$。
  但我们刚刚证明了如果 $\Gamma^* \Proves !B \lor !B$，那么 $!B \in \Gamma^*$。
  因此，$\Gamma^*$ 满足 \olref{defn:prime} 的 \olref{defn:prime2} 条件。
\end{proof}

\begin{prob}
证明：如果在经典逻辑中 $\Gamma \Proves/ \bot$，那么 $\Gamma$ 是一致的，即存在一个赋值使 $\Gamma$ 中的所有公式为真。
\end{prob}

\end{document}